
% ===========================================================================
\section{Averaging as an LMO over Past Samples}
\label{app:matnormal}
\Cref{alg:ours} averages over past batches in two places: the factor
momenta $\widehat M_A,\widehat M_B$ replace the factor gradients in the
direction step~\eqref{eq:factor-directions}, and an EMA accumulates the
vectors $p,q$ that form the curvature factors $P,Q$. Both averages are
standard practice, and this appendix interprets them through the
per-sample LMO~\eqref{eq:sample-lmo} extended over all past samples
with decaying weights.
\Cref{sec:weighted-lmo} states this weighted extension and shows that
the history enters the problem only through two
weighted moments, an average of the batch gradients and a second moment
of the per-sample gradients.
\Cref{sec:weighted-estimation} estimates
both moments from the factor gradients alone, obtaining the $p,q$
recursion of \Cref{alg:ours} and recovering the
update~\eqref{eq:polora-update}. \Cref{sec:scale-placement} explains
why \Cref{alg:ours} normalizes the accumulated vectors $p,q$ after the
average rather than normalizing each per-step estimate before it. The
update is invariant to the scale of the curvature factors, but the
average that estimates them depends on that scale.


\subsection{The Weighted LMO}
\label{sec:weighted-lmo}
A single minibatch gives a noisy objective and only a handful of
constraints, so an update fit to it alone overfits that minibatch. To
counter this, we extend the LMO~\eqref{eq:sample-lmo} over past
minibatches. Ideally every past sample would enter with equal weight,
since all are equally informative once scored at the current weights.
What forces a weighting is that each past gradient is available only
at its own step's weights, and substituting it for the current-weight
gradient errs in proportion to the drift since that step. The decaying
weights suppress the distant steps, where this error is large, and
\Cref{as:slow-drift} below treats the drift within the remaining short
window as negligible. This recency bias is exactly what the momentum
and EMA choices below implement.

\paragraph{The weighted problem.} Let $s$ index the optimizer steps and
$i$ the samples of the size-$n_s$ minibatch at step $s$. Write
$G_{s,i}=\nabla\ell_{x_{s,i}}(W_s)$ for the per-sample gradient, $g_{s,i}=\vec(G_{s,i})$, and
$G_s=\frac1{n_s}\sum_iG_{s,i}$ for the  batch gradient.

Let
$m_{t,s}$ and $w_{t,s}$ be nonnegative weights on the past steps, one
for the objective and one for the constraints. The argument above
fixes only that the weights should decay with age, not how fast,
so we keep them general and defer the choice to~\eqref{eq:EMAweights}
below. We now rewrite the loss and constraint of~\eqref{eq:sample-lmo}
with these weights, both evaluated at the current weights $W_t$:
\begin{equation}\label{eq:weighted-sample-lmo}
 \begin{array}{ll}
  \underset{\Delta W}{\mathrm{minimize}} &
  \displaystyle\sum_{s\le t} \frac{m_{t,s}}{n_s}\sum_{i=1}^{n_s}
    \ip{\nabla\ell_{x_{s,i}}(W_t), \Delta W} \\[2pt]
  \mbox{subject to} &
  \sqrt{w_{t,s}/n_s}\,\bigl|\ip{\nabla\ell_{x_{s,i}}(W_t), \Delta W}\bigr|
  \le \tau, \qquad s=1, \ldots, t,\quad i=1,\dots,n_s.
  \end{array}
\end{equation}
%%\rob{There's a lot of terms/adjectives that make for poor math writing. For instance ..etc Terms that don't mean anything precise, are not defined, and only used once. This is a type of writing that works well in a more literary setting, but only adds confusion in a math/science setting. I suspect this is because some of the text is generated.}
Problem~\eqref{eq:weighted-sample-lmo} uses the gradients
$\nabla\ell_{x_{s,i}}(W_t)$, available only for the current batch
$s=t$. We first apply the outergradient relaxation of
\Cref{sec:sample-loss-lmo}, then substitute the gradients computed at
each past step via \Cref{as:slow-drift}.

\paragraph{The two weighted moments.} Applying the outergradient
relaxation of \Cref{sec:sample-loss-lmo} to~\eqref{eq:weighted-sample-lmo}
leaves a problem that depends on the history only through two
weighted moments, given by
\[
\overline G_t=\sum_{s\le t}\frac{m_{t,s}}{n_s}\sum_{i=1}^{n_s}\nabla\ell_{x_{s,i}}(W_t),
\qquad
\overline\Sigma_t = \sum_{s\le t}\frac{w_{t,s}}{n_s}\sum_{i=1}^{n_s}
  \vec(\nabla\ell_{x_{s,i}}(W_t))\vec(\nabla\ell_{x_{s,i}}(W_t))^\top.
\]
That is, $\overline G_t$ is the weighted average of the batch
gradients, and $\overline\Sigma_t$ is the weighted second moment of
the per-sample gradients. The objective in~\eqref{eq:weighted-sample-lmo}
already equals $\ip{\overline G_t,\Delta W}$ by linearity. For the
constraints, the outer products of the rescaled gradients
sum to $\overline\Sigma_t$, so the leverage bound of \Cref{lem:gradprop}
places each of them in the outergradient set
$\mathcal{G}(\overline\Sigma_t)$ of~\eqref{eq:outgrad} and the
constraints relax exactly as in~\eqref{eq:polorafull-hard}. The relaxed problem is
\begin{equation}\label{eq:weighted-relaxed}
 \begin{array}{ll}
  \underset{\Delta W}{\mathrm{minimize}} &
  \ip{\overline G_t, \Delta W} \\[2pt]
  \mbox{subject to} &
  \displaystyle\max_{\widetilde{G}\in\mathcal{G}(\overline\Sigma_t)}\ip{\widetilde{G},\Delta W}\le\tau .
  \end{array}
\end{equation}

\paragraph{From ideal to computed gradients.} Both moments are built from
the current-weight gradients $\nabla\ell_{x_{s,i}}(W_t)$, unavailable for
past samples since sample $(s,i)$ was backpropagated at its own step
$W_s$. One standing assumption closes this gap, and every later step
invokes it.
\begin{assumption}[Frozen window]\label{as:slow-drift}
Over the effective window of the weights $m_{t,s}$ and $w_{t,s}$, the
drift of the trainable factors is negligible: $A_s\approx A_t$ and
$B_s\approx B_t$ for all steps $s$ in the window.
\end{assumption}
\rob{The awkward thing about this assumption is that it is not precise. What is the "effective window of the weights"? I think it would be best to just state: We replace the gradient and the covariance with momentum style estimates. This attempt at precisely justifying open more questions than it gives answers.}
\noindent As a result, a step-$s$ quantity inside a window average may
be exchanged for its step-$t$ value, with error first order in the
drift $A_t-A_s$, $B_t-B_s$. We use two such exchanges: the
\emph{gradient exchange}, where the computed gradient stands in for
the current-weight gradient,
$G_{s,i}=\nabla\ell_{x_{s,i}}(W_s)\approx\nabla\ell_{x_{s,i}}(W_t)$ by
smoothness of $\ell$, and the \emph{factor exchange}, where factor
expressions keep their current values, as in
$B_s^\top G_s\approx B_t^\top G_s$. The gradient exchange reduces the
moments to
\[
\overline G_t\approx\sum_{s\le t}m_{t,s}G_s,
\qquad
\overline\Sigma_t \approx \sum_{s\le t}\frac{w_{t,s}}{n_s}\sum_{i=1}^{n_s}g_{s,i}g_{s,i}^\top,
\]
computable from the per-sample gradients as they arrive.



\paragraph{Choosing the weights.}
A valid weighting concentrates its mass on the window where the
exchanges of \Cref{as:slow-drift} hold, and a practical one must be
computable without storing the past. Geometric decay meets both
requirements, and the constant-memory recursions below induce it. For
the objective in \eqref{eq:weighted-relaxed}, we set $\overline G_t$ to a
look-ahead momentum $\widehat M_t$ of the batch gradients, mirroring the
momentum scheme of \Cref{alg:ours},
\[
  M_{t+1}=\beta_1M_t+(1-\beta_1)G_t,
  \qquad
  \widehat M_t=\beta_1M_{t+1}+(1-\beta_1)G_t.
\]
This choice emphasizes decreasing the loss on recent batches. For the
constraints in \eqref{eq:weighted-relaxed}, we use an EMA of the within-batch second
moments,
\[
  \overline\Sigma_t
  =\beta_2\,\overline\Sigma_{t-1}
  +(1-\beta_2)\,\frac1{n_t}\sum_{i=1}^{n_t}g_{t,i}g_{t,i}^\top.
\]
This choice gives tighter constraints on recent batches. Unrolling the
recursions gives the induced weights
\begin{equation}\label{eq:EMAweights}
  m_{t,s}=\beta_1^2(1-\beta_1)\beta_1^{t-1-s}
  \quad(s<t),
  \qquad
  m_{t,t}=1-\beta_1^2,
  \qquad
  w_{t,s}=(1-\beta_2)\beta_2^{t-s}.
\end{equation}

\paragraph{Approximating the second moment.} As discussed in~\Cref{sec:sample-loss-lmo},
$\overline\Sigma_t$ cannot be used directly, for two reasons: it is
built from per-sample gradients that backpropagation never
materializes, and it is too large to store. Two approximations remove
these obstacles in turn.
\begin{itemize}[leftmargin=1.5em,itemsep=2pt,topsep=2pt]
  \item \emph{Batch collapse.} We replace the within-batch average
  $\frac1{n_s}\sum_ig_{s,i}g_{s,i}^\top$ by the outer product
  $g_sg_s^\top$ of the batch gradient $g_s=\vec(G_s)$. This substitution
  underlies the second moments of AdaGrad~\cite{adagrad} and its
  relatives such as Adam~\cite{adam} and Shampoo~\cite{shampoo}, which
  all accumulate outer products of batch gradients, and is exact at
  batch size one. Under it the collapsed estimate is a
  weighted average over batch gradients,
  \[
    \overline\Sigma_t\approx\E_{s \leq t}[g_sg_s^\top],
    \qquad\text{where}\qquad
    \E_{s \leq t}[\phi(G_s)]=\sum_{s\le t}w_{t,s}\,\phi(G_s).
  \]
  \item \emph{Kronecker factorization.} We assume the
  Kronecker-factored form of~\eqref{eq:pq-factor} with diagonal
  factors, written $P^\star,Q^\star$ to distinguish the model factors
  from the estimates $P,Q$ that \Cref{alg:ours} stores. Since
  $Q^\star\otimes P^\star$ is unchanged by the rescaling
  $(P^\star,Q^\star)\mapsto(aP^\star,a^{-1}Q^\star)$, we fix
  $\norm{\diag P^\star}_\infty=\norm{\diag Q^\star}_\infty=1$ and let an additional scalar
  $\kappa_t>0$ carry the magnitude:
  \begin{equation}\label{e-diag-kron-second-mom-model}
    \E_{s \leq t}[g_sg_s^\top]\approx\kappa_t\,(Q^\star_t\otimes P^\star_t).
  \end{equation}
\end{itemize}

\subsection{Estimating the Weighted Moments}
\label{sec:weighted-estimation}
In LoRA training, backpropagation gives the factor gradients $G_{A,s}=B_s^\top G_s$
and $G_{B,s}=G_sA_s^\top$~\eqref{e-grad-wrt-factors} at no extra cost. Forming
the full gradient $G_s$, by contrast, would take a $d_{\mathrm{out}}\times
d_{\mathrm{in}}$ matmul and full-matrix storage per layer, exactly what
the low-rank structure avoids. We estimate both weighted moments of \Cref{sec:weighted-lmo}
from the factor gradients alone, taking the momentum $\widehat M_t$ for
the objective and the model factors $P^\star,Q^\star$ for the
constraint.
% , with each replacement of a past-step factor by its current
% value made under \Cref{as:slow-drift}.

\paragraph{Factor momenta.} Under the linearized product
update~\eqref{eq:DeltaW-lin}, the objective of~\eqref{eq:weighted-relaxed}
decomposes as
\[
  \ip{\widehat M_t,\,\Delta B\,A+B\,\Delta A}
  =\ip{\widehat M_tA^\top,\Delta B}+\ip{B^\top\widehat M_t,\Delta A},
\]
so the update depends on $\widehat M_t$ only through
$\widehat M_tA^\top$ and $B^\top\widehat M_t$. The factor momenta
$\widehat M_A,\widehat M_B$ of \Cref{alg:ours} (line~\ref{ln:mom})
estimate precisely these, with $\widehat M_A$ averaging $B_s^\top G_s$
and $\widehat M_B$ averaging $G_sA_s^\top$ under the same weights
$m_{t,s}$. By the factor exchange of \Cref{as:slow-drift},
$\widehat M_A\approx B_t^\top\widehat M_t$ and
$\widehat M_B\approx\widehat M_tA_t^\top$.

\paragraph{Moments of the factor gradients.} To estimate $P^\star$ and $Q^\star$ from the factor gradients alone, we link their second moments to the second moment of the batch gradient. To this end, note that
$\vec(G_{A,s})=(I\otimes B_s^\top)\vec(G_s)$. Consequently the second moment of the gradient of $A$ is given by
\begin{align*}
 \Sigma_{A,t} &:=  \E_{s \leq t}\big[ \vec(G_{A,s})  \vec(G_{A,s})^\top\big]  \\
 &=  \E_{s \leq t}\big[ (I\otimes B_s^\top)\vec(G_s)\vec(G_s)^\top   (I\otimes B_s)\big]  \\
   &\approx  (I\otimes B_t^\top) \E_{s \leq t}\big[ g_s g_s^\top   \big]  (I\otimes B_t),
\end{align*}
\rob{Again \Cref{as:slow-drift} is not precise enough. Here we need to say explicitly, our approximation replaces $B_s$ by $B_t$ for all $s \leq t.$ }
where the last line uses the factor exchange of \Cref{as:slow-drift} to bring the current factor $B_t$ out of the sum. Using~\eqref{e-diag-kron-second-mom-model} in the above gives
\begin{align*}
 \Sigma_{A,t} 
   &\approx  \kappa_t\,  (I\otimes B_t^\top)  (Q^\star_t\otimes P^\star_t)  (I\otimes B_t) \; = \kappa_t\,   Q^\star_t\otimes  (B_t^\top P^\star_t B_t),
\end{align*} 
where the equality uses the mixed-product rule of Kronecker products.
Using analogous arguments for the gradient of $B$, we arrive at the following approximations
\begin{align}
  \Sigma_{A,t} & =\E_{s \leq t}\big[\vec(G_{A,s})\vec(G_{A,s})^\top\big] \approx\kappa_t\,Q^\star_t\otimes C^\star_{A,t} \nonumber \\
  \Sigma_{B,t}&=\E_{s \leq t}\big[\vec(G_{B,s}^\top)\vec(G_{B,s}^\top)^\top\big] \approx \kappa_t\,P^\star_t\otimes C^\star_{B,t}, \label{eq:observable-moments}
\end{align}
where $C^\star_{A,t}=B_t^\top P^\star_tB_t$ and $C^\star_{B,t}=A_tQ^\star_tA_t^\top$ are the $r\times r$ matrices of~\eqref{eq:ours-closure} evaluated at the model factors.

\paragraph{Fitting the factors.} We approximate the model factors $P^\star$
and $Q^\star$ by fitting the second-moment model to the observed
moments~\eqref{eq:observable-moments}. For brevity, here we drop the subscript $t$ since
all quantities are at the current step. Following
KL-Shampoo~\cite{klshampoo}, we fit with the log-determinant
divergence~\cite{logdetdiv}
\[
  D(\Sigma,S)=\Tr(\Sigma S^{-1})-\log\det(\Sigma S^{-1})-d,
\]
defined for positive definite matrices $\Sigma$ and $S$ of common
dimension $d$.
% Away from the exact model, no candidate reproduces the
% observed moment, and the divergence decides how the mismatch is
% weighed. 
Using~\eqref{eq:observable-moments} we fit the two factors by solving
\begin{equation}\label{eq:factor-fit}
  \hat{q}(\widehat P)
  =\argmin_{q>0}\;
  D\big(\Sigma_A,\; \diag(q) \otimes \widehat C_A\big),
  \qquad
  \hat{p}(\widehat Q)
  =\argmin_{p>0}\;
  D\big(\Sigma_B,\;\diag(p) \otimes \widehat C_B\big),
\end{equation}
where we freeze the current estimate of $\widehat C_A$ when estimating $\hat{q}$, and freeze the estimate of $\widehat C_B$ when estimating $\hat{p}$. This alternating freezing is needed to get a closed form update for $\hat{q}$ and $\hat{p}$, since otherwise $\widehat C_A$ and $\widehat C_B$ introduce a joint dependency through
$\widehat C_A=B^\top\widehat PB$ and
$\widehat C_B=A\widehat QA^\top$, for diagonal
$\widehat P,\widehat Q\succ0$.

% Because $C^\star_A$ contains $P^\star$ and $C^\star_B$
% contains $Q^\star$, each fit must be posed with an estimate of the
% other factor, $\widehat C_A=B^\top\widehat PB$ and
% $\widehat C_B=A\widehat QA^\top$ for diagonal
% $\widehat P,\widehat Q\succ0$, which couples the two fits:
% \begin{equation}\label{eq:factor-fit}
%   \hat{q}(\widehat P)
%   =\argmin_{q>0}\;
%   D\big(\Sigma_A,\; \diag(q) \otimes \widehat C_A\big),
%   \qquad
%   \hat{p}(\widehat Q)
%   =\argmin_{p>0}\;
%   D\big(\Sigma_B,\;\diag(p) \otimes \widehat C_B\big).
% \end{equation}
The closed form solutions to~\eqref{eq:factor-fit} are given by
\[
  \hat{q}(\widehat P)=\frac1r\diag\!\left(\E_{s \leq t}[G_{A,s}^\top \widehat C_A^{-1}G_{A,s}]\right),
  \qquad
  \hat{p}(\widehat Q)=\frac1r\diag\!\left(\E_{s \leq t}[G_{B,s} \widehat C_B^{-1}G_{B,s}^\top]\right).
\]
For fixed positive definite $\Sigma$ and $S$, let
$$c^\star(\Sigma,S) :=\argmin_{c>0}D(\Sigma,cS)=\Tr(\Sigma S^{-1})/d$$
be the best scalar multiple
of $S$.
The
next lemma shows the coupled fits are consistent even when posed with
inexact estimates.

\begin{lemma}[Consistency]\label{lem:consistency}
Let the model~\eqref{e-diag-kron-second-mom-model} hold with equality. For any
diagonal $\widehat P\succ0$ and $\widehat Q\succ0$, the
fits~\eqref{eq:factor-fit}  with $\widehat C_A=B^\top\widehat PB$
and $\widehat C_B=A\widehat QA^\top$ return
\[
  \hat{q}(\widehat P)=\kappa\,c^\star\big(C^\star_A,\widehat C_A\big)\,\diag(Q^\star),
  \qquad
  \hat{p}(\widehat Q)=\kappa\,c^\star\big(C^\star_B,\widehat C_B\big)\,\diag(P^\star),
\]
so $\hat{q}/\norm{\hat{q}}_\infty=\diag(Q^\star)$ and
$\hat{p}/\norm{\hat{p}}_\infty=\diag(P^\star)$.
\end{lemma}
\begin{proof}
\rob{This proof also assumes~\eqref{eq:observable-moments} holds to equality.}
Under the model, \eqref{eq:observable-moments} makes $\Sigma_A$ block
diagonal with $j$th block $\kappa\,q^\star_j\,C^\star_A$, and the candidate
$\diag(q)\otimes\widehat C_A$ of~\eqref{eq:factor-fit} at a point $q>0$
is block diagonal with $j$th block $q_j\,\widehat C_A$. Traces and
log-determinants add across blocks, so the objective is
$\sum_jD(\kappa q^\star_jC^\star_A,\,q_j\widehat C_A)$, one term per
block, and the $j$th term is minimized at its best scalar
multiple,
\[
\hat{q}_j=c^\star(\kappa q^\star_jC^\star_A,\widehat C_A)
=\kappa\,q^\star_j\,c^\star(C^\star_A,\widehat C_A).
\]
The factor
$\kappa\,c^\star(C^\star_A,\widehat C_A)$ is the same for every $j$, and
$\norm{\diag Q^\star}_\infty=1$, so the normalization removes it. The claim
for $\hat{p}$ follows with the roles of the two factors exchanged.
\end{proof}

\paragraph{Online estimator.}
The fit~\eqref{eq:factor-fit} is a weighted average over the past factor
gradients with the geometric weights $w_{t,s}$ of
\Cref{sec:weighted-lmo}, so \methodname{} forms it online by the same
EMA. At each step it evaluates the single-batch fit
of~\eqref{eq:factor-fit} on the current batch (write
$\hat{q}_t,\hat{p}_t$), posed with the stored factors
$C_{A,t}=B_t^\top P_tB_t$ and $C_{B,t}=A_tQ_tA_t^\top$, and folds the
result into a running average,
\begin{equation}\label{eq:ours-klcoupling}
\begin{aligned}
  \hat{q}_t(P_t) &= \frac1r\diag\!\left(G_{A,t}^\top C_{A,t}^{-1}G_{A,t}\right),
  &
  \hat{p}_t(Q_t) &= \frac1r\diag\!\left(G_{B,t} C_{B,t}^{-1}G_{B,t}^\top\right),
  \\
  q_{t+1} &= \beta_2q_t+(1-\beta_2)\hat{q}_t(P_t),
  &
  p_{t+1} &= \beta_2p_t+(1-\beta_2)\hat{p}_t(Q_t),
  \\
  Q_{t+1} &= \diag\!\bigl(q_{t+1}/\norm{q_{t+1}}_\infty\bigr),
  &
  P_{t+1} &= \diag\!\bigl(p_{t+1}/\norm{p_{t+1}}_\infty\bigr).
\end{aligned}
\end{equation}
The last line applies the normalization that \Cref{lem:consistency}
justifies. Unrolled, the recursion weights the per-step estimates by
exactly the $w_{t,s}$ of \Cref{sec:weighted-lmo}. Within the window,
$B_s\approx B_t$ by \Cref{as:slow-drift}, and
the stored $P_s$ changes only at rate $1-\beta_2$, so
$C_{A,s}\approx C_{A,t}$ and each per-step estimate already uses the
current matrices. The running $q_t$ then matches the full fit
$\hat{q}(P_t)$, the consistent fit of \Cref{lem:consistency}.

\paragraph{Comparison with KL-Shampoo.} KL-Shampoo~\cite{klshampoo}
also fits Kronecker factors by minimizing the log-determinant
divergence, and its moving average has the same shape
as~\eqref{eq:ours-klcoupling}. Our approach differs in the following ways:
\begin{itemize}[leftmargin=1.5em,itemsep=1pt,topsep=2pt]
  \item \emph{No distributional assumptions.} KL-Shampoo fits the second moment
  of a stationary gradient distribution. Here $\E_{s \leq t}$ is a weighted sum
  over the realized gradients, not an expectation, and the statements
  above hold for an arbitrary gradient stream.
  \item \emph{The observables are the factor gradients.} KL-Shampoo fits
  dense factors to the second moment of the weight gradient $G$. Here
  each diagonal factor is fit to the second moment of its own factor
  gradient, $G$ is never formed, and the two fits are coupled through
  $C_A$ and $C_B$.
  \item \emph{Valid at every step.} KL-Shampoo justifies its
  moving average as a stochastic proximal-gradient step on the
  divergence objective, so its estimates are correct in the limit. Here
  the recursion computes the weighted moments exactly at every step,
  and \Cref{lem:consistency} shows that posing a fit with an inexact
  estimate of the other factor only rescales its output, a scale the
  normalization removes. The one remaining approximation is the frozen
  window of \Cref{as:slow-drift}.
\end{itemize}

\paragraph{Recovering the update of \Cref{alg:ours}.} Substituting the
fitted model $\kappa_t(Q_t\otimes P_t)$ for $\overline\Sigma_t$ and the
look-ahead momentum $\widehat M_t$ for $\overline G_t$ turns the
weighted problem~\eqref{eq:weighted-relaxed} into the preconditioned
spectral LMO~\eqref{eq:polorafull}
(\Cref{lem:pseudogradient-spectral}). Solving it with the direction and
magnitude steps of \Cref{sec:sample-loss-lmo}, with the factor momenta
standing in for the factor gradients, returns the
update~\eqref{eq:polora-update}. The only departure of \Cref{alg:ours}
from this derivation is that the curvature factors update after the
direction step, so step $t$ runs on the estimate from step $t-1$, which
avoids a second formation and inversion of $C_A$ and $C_B$ per step.

\subsection{Normalization Placement}
\label{sec:scale-placement}
\methodname{} normalizes the accumulated vectors $q_{t+1},p_{t+1}$ to unit
largest entry when it forms the next-step factors
$Q_{t+1},P_{t+1}$~\eqref{eq:ours-klcoupling}, rather than normalizing each
per-step estimate $\hat{q}_t,\hat{p}_t$ before it enters the average.
Normalizing the accumulated vectors does not change the update, which
depends only on the direction of the stored factors. Normalizing each
estimate first would change the average, which is sensitive to the scale
of those estimates.

\paragraph{The update is scale-free.} The update depends on the stored
factors only through their direction. Under a rescaling $Q\mapsto cQ$
($c>0$), both $Q^{-1/2}$ and $C_B^{-1/2}=(AQA^\top)^{-1/2}$ pick up a
factor $c^{-1/2}$. Inside the matrix sign this scalar has no effect, since
the sign is scale-invariant. Outside it, the scalar multiplies $D_A$ and
$D_B$ by $c^{-1/2}$, which the division by $\norm{D_A}_2$ and
$\norm{D_B}_2$ in the magnitude step cancels. The factor $P$ enters the
same way. The update therefore sees only the direction of the stored
factors, never the magnitude $\kappa_t$ of the
model~\eqref{e-diag-kron-second-mom-model}, so normalizing the accumulated
vectors leaves it unchanged.

\paragraph{The average is scale-sensitive.} The coupled fit makes each
per-step estimate scale inversely with the stored factor it is computed
from,
\[
  \hat{q}_t(aP_t)=a^{-1}\,\hat{q}_t(P_t),
  \qquad
  \hat{p}_t(aQ_t)=a^{-1}\,\hat{p}_t(Q_t)
  \qquad(a>0),
\]
since $C_A=B^\top PB$~\eqref{eq:ours-closure} is linear in $P$, and
likewise $C_B$ in $Q$. Normalizing after
averaging gives every stored factor the same scale
($\norm{\diag P_t}_\infty=1$ at every step), so the estimates enter the
average on a common footing. Were the stored factor at step $s$ left carrying a step-varying scale
$a_s$ instead, the identity above would scale
$\hat{q}_s$ by $1/a_s$, so it would enter with effective weight
$w_{t,s}/a_s$ rather than the $w_{t,s}$ chosen in \Cref{sec:weighted-lmo}.
Normalizing the per-step estimates themselves before averaging would fail
for a second reason. Normalization is nonlinear and does not commute with
the mean, so per-step noise that cancels in the raw average would survive.
This average-then-normalize order is the one by which momentum precedes
normalization in normalized SGD~\cite{cutkosky_nigt} and orthogonalization
in Muon~\cite{denoise_ortho}.

\section{Numerical Subroutines}
\label{app:primitives}

\subsection{Spectral-Norm Estimation}
\label{app:smax}
\Cref{alg:ours} requires repeated estimates of spectral norms.  
We estimate each $\norm{M}_2$ by warm-started power iteration on the smaller of
the two Gram matrices (\Cref{alg:smax}), caching the leading vector across
optimizer steps so it starts near the top singular vector.  A cold or stale start
can under-estimate $\norm{M}_2$; since the optimizer divides by this estimate,
that would inflate the update and risk destabilizing training.  We therefore
return the larger of the power-iteration value and the largest row norm of $M$, a
lower bound on $\norm{M}_2$.  We use $K=8$ iterations.

\begin{algorithm}[ht]
\caption{Guarded spectral-norm estimate of $\norm{M}_2$}
\label{alg:smax}
\begin{algorithmic}[1]
\Require $M\in\mathbb{R}^{m\times n}$, $m\le n$ (else $M^\top$); cached $v\in\mathbb{R}^m$; iterations $K$
\State $l \gets \max_i\norm{M_{i,:}}_2$
\If{$v$ is empty} \State $v\gets M\mathbf{1}$ \Comment{cold start} \EndIf
\For{$k=1,\dots,K$} \Comment{iterate the smaller Gram $MM^\top$}
  \State $w\gets M^\top v$;\quad $v\gets Mw/\norm{Mw}_2$
\EndFor
\State \Return $\max\big\{\norm{M^\top v}_2,\ l\big\}$
\end{algorithmic}
\end{algorithm}